\subsection{$(B,A,A)$ 型膜作为工具}\label{lagbr}

$(B,A,A)$ 型膜是关于复结构 $I$ 的 $B$-膜，又是关于 $\omega_J$ 与 $\omega_K$ 的 $A$-膜。它们在第 \ref{genera} 节引入。我们特别关心秩 1 的 $(B,A,A)$ 型拉格朗日膜。这样膜的支撑是复拉格朗日子流形 $L$。我们已取 $B$-场为 $\sB=\omega_J$，它限制到复拉格朗日子流形 $L$ 上为零。因此支撑在 $L$ 上的秩 1 $A$-膜 $\B_0$ 的 $\CP$ 曲率条件 $\sF+\sB=0$ 化为 $\sF=0$。故这样膜的 $\CP$ 丛是平坦 $\spinc$ 结构。如第 \ref{genera} 节所述，若 $\B_0$ 是这样的膜，则 $\Hom(\B_0,\B_\cc)$ 是 $Y$ 上函数量子形变代数的模，尽管这个模不能自然地解释为 $L$ 的量子化。

我们可以用这样的模作为理解 $\Hom(\B_\cc,\B_\cc)$ 的工具。作为探针，考虑一个方便的 $(B,A,A)$ 型膜，即支撑在余切丛零截面 $W\subset T^*W$ 上的膜。我们假定 $W$ 容许平坦 $\spinc$ 结构 $\L$，因而是一个秩 1 拉格朗日膜 $\B_0$ 的支撑。由于考虑的是沿 $W$ 局部的论证，$\L$ 是否整体存在其实不重要；同样，稍后我们会看到 $\L$ 的选取也无关紧要。

余切丛的零截面在纤维伸缩下不变。故 $\Hom(\B_0,\B_\cc)$ 可在 $\hbar$ 很小时计算。考虑到 $\B_\cc$ 的 $\CP$ 丛平凡，按附录 \ref{details} 的一般分析，结果是 $\Hom(\B_0,\B_\cc)\cong\oplus_{q=0}^{\dim_\C W}H^q(W,K_W^{1/2}\otimes\L^{-1})$。特别地，只取 $q=0$，$\Hom(\B_\cc,\B_\cc)$ 必须能作用在 $H^0(W,K_W^{1/2}\otimes\L^{-1})$ 上，即 $K_W^{1/2}\otimes\L^{-1}$ 的全纯截面空间。由于 $\L^{-1}$ 平坦，它在此不起作用：作用在 $K_W^{1/2}\otimes\L^{-1}$ 截面上的全纯微分算符，与作用在 $K_W^{1/2}$ 截面上的全纯微分算符相同。因此 $\Hom(\B_\cc,\B_\cc)$ 必须是作用在 $K_W^{1/2}$ 截面上（而不是某个其它线丛上）的微分算符层。

按先前解释，作用在 $K_W^{1/2}$ 上的微分算符无论 $K_W^{1/2}$ 是否整体作为线丛存在都有意义。不过或许值得说明：类似地，若 $K_W^{1/2}$ 确实作为线丛存在，它的选取也无关紧要。$K_W^{1/2}$ 的不同选取相差一个满足 $\RR^2$ 平凡的线丛 $\RR$ 的张量积。这样的 $\RR$ 容许平坦联络，故张量以 $\RR$ 对作用在 $K_W^{1/2}$ 上的全纯微分算符没有影响。

在附录 \ref{opposite} 中，我们解释另一种（数学上更标准的）办法，说明在余切丛情形 $\Hom(\B_\cc,\B_\cc)$ 是作用在 $K_W^{1/2}$ 截面上的微分算符层，而不是线丛的某个其它复幂。
